% index_parameters.tex
\documentclass[12pt]{article}
\usepackage[utf8]{inputenc}
\usepackage[paper=a4paper,margin=2.5cm]{geometry}
\begin{document}

\section*{אינדקס parameters.tex}
\textbf{מטרה}:
להגדיר ולהסביר כל פרמטר שנמצא בפרומט המחקרי.

\subsection*{פרמטרים מרכזיים ותיאור קצר}
\begin{itemize}
  \item \textbf{Title} שם קצר ותמציתי.
  \item \textbf{Objective} ניסוח מטרה מדויק.
  \item \textbf{Scope} טווח זמן, אוכלוסייה, תחום.
  \item \textbf{Research Questions} שאלות מחקר מנוסחות.
  \item \textbf{Methodology} סוגי שיטות ופורמט תיאורן.
  \item \textbf{Data Sources} סוגי מקורות והנחיות לאימות.
  \item \textbf{Inclusion/Exclusion} קריטריונים מפורטים.
  \item \textbf{Extraction Fields} שדות חילוץ סטנדרטיים.
  \item \textbf{Deliverables} פורמטים רצויים לתוצרים.
  \item \textbf{Constraints} סוגי מגבלות ודרישות אתיות.
  \item \textbf{Evaluation Metrics} מדדי איכות וכיצד לחשבם.
  \item \textbf{Tone and Style} דוגמאות משפטיות וסגנוניות.
  \item \textbf{Length and Format} פורמטים נתמכים: LaTeX, Markdown, JSON.
\end{itemize}

\subsection*{הפניות}
עבור דוגמאות ערכים ראו \textbf{index\_prompt\_generator.tex} ו-\textbf{index\_README.tex}.
\end{document}
